\section{条件认证:有限样本上界}\label{sec:cert}

目标:在模型类内,\emph{直接}对大码距的 $\pL(d)$ 给出有限样本上界,而不是先估计 $\alpha$ 再外推。
证明骨架分三层:\emph{(A) 统计层}(由数据得到参数的置信上界);\emph{(B) 单调层}(参数的上界给出模型的``支配'',
由定理~\ref{thm:monotone} 得 $\err$ 的上界);\emph{(C) 计算层}(对支配模型给出 $\err_d$ 的显式上界)。
本节完整证明码容量情形(定理~\ref{thm:cc-cert}),并把含突发与芯片事件的情形(定理~\ref{thm:burst-cert})
化为三个明确的开放引理。

\paragraph{关于“哪个解码器”。}
本节所有上界都是对\emph{最大似然(Bayes 最优)解码}的上界,表述为“存在解码器(已知噪声参数下的 ML)
使逻辑失败概率不超过 $U$”。由最优性,它\emph{不是}对实验中实际使用的次优解码器(如 MWPM)失败率的上界。
反之,本节的\emph{下界}(地板,引理~\ref{lem:floor})对任何解码器成立。

\subsection{(C) 联合(Bhattacharyya)上界}

\begin{lemma}[联合上界]\label{lem:union}
设 $(E,F)$ 为线性解码结构 $(V,\sigma,L)$ 上的任意噪声实例,令 $K=\ker\sigma\setminus\ker L$
(不被检测却翻转可观测量的故障)。则
\[
\err(E,F)\;\le\;\tfrac12\sum_{G\in K}\BC(P;G),\qquad
\BC(P;G):=\sum_{E,F}\sqrt{P(E,F)\,P(E+G,F)}.
\]
\hfill\st{PROVED}
\end{lemma}
\begin{proof}
固定观测 $y=(s,f)$,记 $P(\ell,y)=\Prob[L=\ell,Y=y]$,ML 判决为 $\ell^\ast(y)$。则
$\err=\sum_y\sum_{\ell\ne\ell^\ast}P(\ell,y)$。对 $\ell\ne\ell^\ast$ 有 $P(\ell,y)\le P(\ell^\ast,y)$,故
$P(\ell,y)\le\sqrt{P(\ell,y)P(\ell^\ast,y)}$,于是
\[
\err\le\sum_y\sum_{\{\ell,\ell'\}}\sqrt{P(\ell,y)P(\ell',y)}\quad(\text{对无序标签对求和}).
\]
记 $c(s,\ell)=\{E:\sigma(E)=s,L(E)=\ell\}$,有 $P(\ell,s,f)=\sum_{A\in c(s,\ell)}P(A,f)$。
由 $\sqrt{\sum_ix_i}\le\sum_i\sqrt{x_i}$ 得
\[
\sqrt{P(\ell,s,f)P(\ell',s,f)}\le\sum_{A\in c(s,\ell)}\sum_{B\in c(s,\ell')}\sqrt{P(A,f)P(B,f)}.
\]
对 $s,\{\ell,\ell'\},f$ 求和,右边变成对满足 $\sigma(A)=\sigma(B)$、$L(A)\ne L(B)$ 的\emph{无序}故障对 $\{A,B\}$ 求和。
令 $G=A+B\in K$,$(A,B)\mapsto(A,G)$ 是有序对到 $V\times K$ 的双射,无序对数为其一半,
故得 $\tfrac12\sum_{G\in K}\sum_{A,f}\sqrt{P(A,f)P(A+G,f)}$。
\end{proof}

\begin{corollary}[码容量的积形式与重量枚举]\label{cor:product}
对 $\Nn_{\rm cc}(p,e)$ 与稳定子码 $(n,S)$,$K=N(S)\setminus S$,且对 $G\in K$
$\BC(P;G)=B(p,e)^{\wt(G)}$,其中 $B=e+(1-e)\beta_{\rm dep}(p)$($\wt$ 为泡利重量)。故
\[
\err_d(p,e)\le\tfrac12\widetilde W_d\big(B(p,e)\big),\qquad
\widetilde W_d(x)=\sum_{G\in N(S)\setminus S}x^{\wt(G)}=W_{N(S)}(x)-W_S(x),
\]
并由 MacWilliams 恒等式 $W_{N(S)}(x)=\dfrac1{|S|}\sum_{g\in S}(1+3x)^{n-\wt(g)}(1-x)^{\wt(g)}$ 可由稳定子群的重量分布计算。
\hfill\st{PROVED}
\end{corollary}
\begin{proof}
单比特:标记 $f=1$(概率 $e$)时分布均匀,平移不变,贡献 $e\cdot1$;标记 $f=0$ 时贡献
$(1-e)\sum_a\sqrt{\Dep_p(a)\Dep_p(a+G_j)}=(1-e)\beta_{\rm dep}(p)$($G_j\in\{X,Y,Z\}$ 三者对称;
配对 $(I,G_j),(G_j,I)$ 给出 $2\sqrt{(1-p)p/3}$,另外两个非平凡泡利互相配对给出 $2p/3$)。
$G_j=I$ 的比特贡献 $1$。独立性给出积形式。MacWilliams 恒等式为标准结果。
\end{proof}

\noindent 对 $d=3$ 的三个码,不等式 $\err\le\tfrac12\widetilde W(B)$ 在 $\texttt{theory/checks}$ 中全部通过(检验 C5)。
一般 $d$ 时 $W_S$ 可由张量网络计算(每格点 4 维腿,转移矩阵宽度随 $d$ 指数增长),具体可行的 $d$ 范围待 M1 之后基准测试,\emph{此处不作承诺}。

\subsection{(A) 统计层:参数的置信上界}

\begin{lemma}[由标记估计 $e$]\label{lem:flags}
$M$ 个独立样本中共 $nM$ 个比特-样本,标记 $F_{ij}$ i.i.d.\ Bernoulli($e$)。令
$\bar e=\bar e(\delta_1)$ 为 Clopper--Pearson 单侧上限,则 $\Prob[e\le\bar e]\ge1-\delta_1$。\hfill\st{PROVED}
\end{lemma}

\begin{lemma}[由综合征频率估计 $\lambda$]\label{lem:lambda}
取一个固定的 $X$ 型稳定子生成元 $g_0$,重量 $w_0$。对其支撑上没有任何标记的那些样本(记其个数为 $M'$),
综合征比特为 i.i.d.\ Bernoulli,参数 $\tfrac12(1-\lambda^{w_0})$(以标记模式为条件,精确成立)。
令 $\overline{P}$ 为该参数的 Clopper--Pearson 上限(置信水平 $1-\delta_2$),则
$\lambda\ge\underline\lambda:=(1-2\overline P)_+^{1/w_0}$,即 $p\le\bar p:=\tfrac34(1-\underline\lambda)$,
概率 $\ge1-\delta_2$。\hfill\st{PROVED}
\end{lemma}
\begin{proof}
没有标记时 $g_0$ 支撑上每个比特以概率 $p$ 去极化;$X$ 型检测器翻转当且仅当支撑上 $\{Y,Z\}$ 型错误的个数为奇数,
单比特贡献本征值 $1-\tfrac43p=\lambda$,故 $\Pr[\text{bit}=1]=\tfrac12(1-\lambda^{w_0})$。
以标记为条件时各样本独立,故为精确二项。
\end{proof}

\begin{theorem}[码容量情形的有限样本认证]\label{thm:cc-cert}
设数据为码容量噪声 $\Nn_{\rm cc}(p,e)$ 的 $M$ 个独立样本(标记与综合征),$\delta=\delta_1+\delta_2$。
令 $\bar e,\bar p$ 如引理~\ref{lem:flags}、\ref{lem:lambda}。则以概率至少 $1-\delta$,对\emph{一切} $d$,
\[
\err_d(p,e)\;\le\;\err_d(\bar p,\bar e)\;\le\;\tfrac12\widetilde W_d\big(B(\bar p,\bar e)\big).
\]
\hfill\st{PROVED}
\end{theorem}
\begin{proof}
在事件 $\{e\le\bar e\}\cap\{p\le\bar p\}$ 上,由定理~\ref{thm:preorder}(情形 $e'\ge e$、$p'\ge p$)
$\Nn_{\rm cc}(\bar p,\bar e)$ 可由 $\Nn_{\rm cc}(p,e)$ 经局域自由操作得到(叠加擦除,再叠加去极化),
故定理~\ref{thm:monotone} 给出第一个不等式;第二个是推论~\ref{cor:product}。
该事件的概率由并集界 $\ge1-\delta_1-\delta_2$。
\end{proof}

这是一个\emph{不依赖任何 $\pL$ 数据}的认证:标记和综合征统计足以给出所有 $d$ 的上界。代价是联合上界的松弛
(它的阈值远低于真实 ML 阈值)。把联合上界换成更紧的、同样单调的上界(例如用测得的小 $d$ 失败率校准)是下一步。

\subsection{含局域突发与芯片尺度事件的情形}

设模型类为定义~\ref{def:events}。记 $q_c=1-e^{-r_c\,T_{\rm tot}}$ 为总时长 $T_{\rm tot}$ 内至少发生一次芯片尺度事件的概率。

\begin{lemma}[芯片尺度事件的地板]\label{lem:floor}
令 ${\err}^{(0)}$ 为去掉芯片尺度层后的同一噪声的 $\err$。则
\[
\err\;\le\;q_c+(1-q_c)\,{\err}^{(0)}.
\]
若进一步假设 \textbf{(K)}:给定发生芯片事件,可观测量 $L(E)\in\F^q$ 均匀且与 $Y=(\sigma(E),F)$ 独立,则
\[
\err\;\ge\;q_c\,(1-2^{-q}).
\]
于是 $\err\sim q_c\propto r_cT_{\rm tot}$ 的“地板”与 $d$ 无关;上界对 (K) 不需要,下界对任何解码器的失败率同样成立。
\hfill\st{PROVED}
\end{lemma}
\begin{proof}
上界:用去掉芯片层后最优的 ML 解码器 $\hat g_0$ 处理完整问题。设 $C$ 为“至少一个芯片事件”。
$\Prob[\hat g_0\text{ 失败}]\le\Prob[C]+\Prob[C^c]\,{\err}^{(0)}$,因为在 $C^c$ 上数据服从去掉芯片层的噪声(各层独立)。
$\err$ 不超过任何解码器的失败率。

下界:考虑知道 $C$ 的先知解码器,其失败率不超过 $\err$(定理~\ref{thm:monotone} 的 (F1) 情形:真实观测是先知观测的粗粒化)。
在 $C$ 上由 (K) 任何猜测的失败概率为 $1-2^{-q}$。
(K) 对“全芯片完全去极化事件”成立:均匀的 $E_{\rm chip}$ 使 $(\sigma,L)(E)$ 在 $\F^{n-k}\times\F^{2k}$ 上均匀(满射),
故 $L$ 与 $\sigma$ 独立;标记独立于 $E_{\rm chip}$。
\end{proof}

\begin{lemma}[突发混合的联合上界]\label{lem:burst-union}
设 $E=E_{\rm bg}+E_{\rm burst}$,以突发构型 $\mathcal C$ 为条件时噪声是比特间独立的积测度,构型分布为 $\pi$。则
\[
\BC(P;G)\;\le\;\sum_{\mathcal C,\mathcal C'}\sqrt{\pi(\mathcal C)\pi(\mathcal C')}\;\prod_{j}\beta_j(\mathcal C,\mathcal C';G_j),
\]
其中 $\beta_j(\mathcal C,\mathcal C';G_j)=\sum_{a,f}\sqrt{P_{\mathcal C,j}(a,f)P_{\mathcal C',j}(a+G_j,f)}$。
\hfill\st{PROVED}
\end{lemma}
\begin{proof}
$P=\sum_{\mathcal C}\pi(\mathcal C)P_{\mathcal C}$。由 $\sqrt{\sum x_i}\le\sum\sqrt{x_i}$,
\[
\sqrt{P(E,F)P(E+G,F)}\le\sum_{\mathcal C,\mathcal C'}\sqrt{\pi(\mathcal C)\pi(\mathcal C')}\sqrt{P_{\mathcal C}(E,F)P_{\mathcal C'}(E+G,F)}.
\]
对 $(E,F)$ 求和并用积测度。
\end{proof}
\noindent 注意不能简单地“以 $\mathcal C$ 为条件、把 $\mathcal C$ 当作已知”:那样得到的是知道构型的解码器的失败率,
它\emph{低于} $\err$(先知更强),不是上界。引理~\ref{lem:burst-union} 是对未知构型的正确处理,但右端把两条独立的突发构型同时求和,
结果偏松。

\begin{theorem}[含突发与芯片事件的条件认证(目标陈述)]\label{thm:burst-cert}
设真实参数 $\theta=(p,e,r,\vartheta,r_c)$,形状 $\vartheta=(R,T_b,p_b)$ 属于已知的有限类 $\Theta_\vartheta$(若问题~\ref{prob:O1} 的单调性成立,可代以形状上限;否则下式右端对 $\vartheta\in\Theta_\vartheta$ 取最大值)。
假设由数据得到的置信上限满足 $\Prob[p\le\bar p,\ e\le\bar e,\ r\le\bar r,\ r_c\le\bar r_c]\ge1-\delta$。则
以概率 $\ge1-\delta$,对一切 $d,T_{\rm tot}$,
\[
\err_{d,T_{\rm tot}}(\theta)\;\le\;\bar q_c+(1-\bar q_c)\,\tfrac12\sum_{G\in K}\ \sum_{\mathcal C,\mathcal C'}\sqrt{\pi_{\bar\theta}(\mathcal C)\pi_{\bar\theta}(\mathcal C')}\prod_j\beta_j(\mathcal C,\mathcal C';G_j),
\]
其中 $\bar q_c=1-e^{-\bar r_cT_{\rm tot}}$,$\pi_{\bar\theta}$ 为参数上限对应的突发构型分布。\hfill\st{PLAN}
\end{theorem}

\noindent\textbf{已证明的组成部分:}速率方向的单调性(命题~\ref{prop:rate-mono}),地板(引理~\ref{lem:floor}),
联合上界(引理~\ref{lem:union}、\ref{lem:burst-union});把 $\bar\theta$ 代入的论证与定理~\ref{thm:cc-cert} 相同。
\textbf{尚缺三块:}
\begin{problem}[O1:形状参数的单调性]\label{prob:O1}
无标记突发下,$\err$ 对 $p_b,R,T_b$ 是否单调?(注~\ref{rem:shape} 说明自由操作不给出;若不成立,
则须对形状取有限类上的最坏情形,或对形状作联合上界。)
\end{problem}
\begin{problem}[O2:污染下的参数识别]\label{prob:O2}
突发与芯片事件污染综合征统计,引理~\ref{lem:lambda} 的 i.i.d.\ 前提失效。需要对 $p$ 的稳健置信上限(例如排除高权重检测事件窗口后的估计)
与对 $r,r_c$ 的置信上限(由事件检测统计,与 P2 的泊松下限 §\ref{sec:design} 衔接)。
\end{problem}
\begin{problem}[O3:突发和的显式估计]\label{prob:O3}
把引理~\ref{lem:burst-union} 右端的二重构型和化为显式的指数函数(团簇展开);需给出收敛域,
并说明在 $d\gg R$ 时它只改变指数、不产生与 $d$ 无关的地板。
\end{problem}

\paragraph{数值检验的分工(见 §\ref{sec:map})。}
电路级 stim + MWPM 的 P2 数据能检验:(i) 地板 $\err\ge q_c(1-2^{-q})$ 与 $\propto r_cT_{\rm tot}$ 的标度
(下界对任何解码器成立,MWPM 的失败率不得低于它);(ii) 参数估计器的无偏性(M3 验收项);
(iii) 突发使指数变小而不是产生地板。\emph{上界}部分针对 ML,MWPM 数据不能证伪它;
须用小 $d$ 的精确 ML 或码容量下带突发簇的张量网络 ML(M1)检验。
